% ==========================================================================
% ATFD: Agent Trajectory Failure Detection — A Benchmark for Evaluating
% Agent Monitoring Tools
% ==========================================================================
% Format: clean article class with 1-inch margins (NeurIPS-compatible)
% ==========================================================================

\documentclass[11pt]{article}

% -- Page geometry --
\usepackage[margin=1in]{geometry}

% -- Fonts & encoding --
\usepackage[utf8]{inputenc}
\usepackage[T1]{fontenc}
\usepackage{times}

% -- Math --
\usepackage{amsmath,amssymb,amsfonts}

% -- Tables & figures --
\usepackage{booktabs}
\usepackage{multirow}
\usepackage{graphicx}
\usepackage{subcaption}
\usepackage{array}
\usepackage{tabularx}
\usepackage{colortbl}

% -- Colors & links --
\usepackage[dvipsnames]{xcolor}
\usepackage[colorlinks=true,linkcolor=blue,citecolor=blue,urlcolor=blue]{hyperref}

% -- Citations --
\usepackage{natbib}
\bibliographystyle{abbrvnat}
\setcitestyle{authoryear,round}

% -- Line numbers (for review) --
\usepackage{lineno}
\linenumbers

% -- Misc --
\usepackage{enumitem}
\usepackage{microtype}
\usepackage{xspace}
\usepackage{algorithm}
\usepackage{algpseudocode}
\usepackage{listings}

% -- Custom commands --
\newcommand{\atfd}{\textsc{ATFD}\xspace}
\newcommand{\taubench}{$\tau$-bench\xspace}
\newcommand{\swebench}{\textsc{SWE-bench}\xspace}
\newcommand{\passk}{pass$^k$\xspace}
\newcommand{\eg}{\emph{e.g.}\xspace}
\newcommand{\ie}{\emph{i.e.}\xspace}
\newcommand{\cf}{\emph{cf.}\xspace}
\newcommand{\etal}{\emph{et al.}\xspace}
\newcommand{\placeholder}[1]{\textcolor{red}{\textbf{#1}}}

% -- Listings style for trajectory examples --
\lstset{
  basicstyle=\ttfamily\footnotesize,
  breaklines=true,
  frame=single,
  numbers=left,
  numberstyle=\tiny\color{gray},
  keywordstyle=\color{blue},
  commentstyle=\color{OliveGreen},
  stringstyle=\color{BrickRed},
}

% ==========================================================================
\title{
  \textbf{ATFD: Agent Trajectory Failure Detection} \\[6pt]
  \large A Benchmark for Evaluating Agent Monitoring Tools
}

\author{
  Anonymous Authors \\
  \texttt{anonymous@example.com}
}

\date{}

\begin{document}
\maketitle

% ==========================================================================
% ABSTRACT
% ==========================================================================
\begin{abstract}
Agent workflows built on large language models are increasingly deployed in
business-critical settings, yet no benchmark exists for evaluating the
\emph{monitoring tools} that must detect failures in these workflows.
Existing benchmarks evaluate agent \emph{capability}---whether an agent
completes a task---but not whether an external system can reliably
\emph{detect} when an agent has failed, partially succeeded, or produced
substandard output. We introduce \atfd (Agent Trajectory Failure Detection),
a benchmark that fills this gap. \atfd contributes: (1)~a 7-category,
23-subcategory failure taxonomy that extends prior work by including
quality degradation as a first-class failure mode; (2)~a formal task
definition where systems must analyze complete agent trajectories and
produce structured findings with severity, category, and attribution;
(3)~a multi-domain dataset of {1,162} trajectories drawn from
\taubench and hand-crafted synthetic scenarios, with \swebench integration
as planned future work; (4)~a
multi-source ground truth protocol combining programmatic verdicts with
LLM-judge consensus; (5)~a comparative evaluation of 7 systems---a
naive heuristic baseline, a zero-config structural heuristic (Galea),
two open-weight LLM judges, Claude Sonnet~4 as a proprietary LLM
judge, LangSmith with custom evaluator rules, and Braintrust with
custom scorers---revealing a stark configuration effort gap; and
(6)~an open benchmark harness with cost-aware metrics. Our evaluation reveals
a stark \emph{detection gap} with complementary profiles: the Galea heuristic
achieves {100\%} detection rate on real \taubench failures with {0\%} false
positives using zero-config pattern matching, but only {78\%} on synthetic
failures that require semantic analysis; conversely, open-weight LLM judges
achieve {100\%} detection on synthetic failures but only {25.0\%} on real
\taubench failures involving subtle argument and state errors.
Claude Sonnet~4 bridges this gap, achieving {100\%} detection on both
synthetic and real retail trajectories with the highest category alignment
(CA = 0.489) of any evaluated system.
This complementarity suggests that ensemble approaches combining structural
and semantic detection could approach near-perfect coverage.
\end{abstract}

% ==========================================================================
% 1. INTRODUCTION
% ==========================================================================
\section{Introduction}
\label{sec:introduction}

Large language model (LLM) agents are transitioning from research prototypes
to production systems that execute consequential business workflows:
customer service interactions, software engineering tasks, financial
analysis, and legal due diligence~\citep{zenml2025production}. This
transition brings an urgent need: when an agent fails---by executing the
wrong action, producing incomplete analysis, or violating an operational
policy---someone or something must detect the failure before it propagates.

The observability landscape has responded. Commercial platforms such as
LangSmith, Langfuse, Arize Phoenix, and Braintrust now offer trace
visualization, LLM-as-judge evaluation, and anomaly detection for agent
workflows. Academic work has proposed moving beyond black-box benchmarking
toward runtime analytics of agentic systems~\citep{moshkovich2025beyond}.
Yet a fundamental question remains unanswered: \emph{how well do these
monitoring tools actually work?}

Today there is no principled way to answer this question. Existing benchmarks
evaluate agent \emph{capability}---whether GPT-4 can resolve a GitHub issue
\citep{jimenez2024swebench}, complete a customer service
task~\citep{yao2025taubench}, or navigate a web
environment~\citep{zhou2024webarena}. These benchmarks produce agent
trajectories and ground-truth outcomes, but they do not evaluate whether a
monitoring tool can correctly detect failures \emph{within} those
trajectories. The distinction matters: a monitoring tool must not only
determine that something went wrong but must also identify \emph{what} went
wrong, \emph{where} in the trajectory it occurred, and \emph{how severe}
the failure is.

The gap is consequential. Production deployments of multi-agent systems
report failure rates between 41\% and
87\%~\citep{zenml2025production,cemri2025mast}. AI safety incidents increased
56.4\% from 2023 to 2024~\citep{responsibleai2025incidents}. Enterprise
adoption demands reliability, but reliability requires monitoring, and
monitoring tools have no benchmark against which to be measured.

\subsection*{Why This Benchmark Is Necessary}

The cost of undetected agent failures is not hypothetical. In July 2025,
Replit's AI coding agent---operating during a declared code freeze---executed
a \texttt{DROP TABLE} command that destroyed months of production data for
over 1,200 users, then fabricated approximately 4,000 fake records in an
apparent attempt to conceal the destruction.\footnote{Fortune, July 2025.
``AI coding tool Replit wiped database, called it a catastrophic failure.''}
A monitoring tool watching for destructive database operations during a
declared freeze period would have flagged this trajectory before the
irreversible action executed.

In the legal domain, attorneys submitted court filings containing entirely
fabricated case citations generated by ChatGPT (\emph{Mata v. Avianca},
S.D.N.Y., 2023), resulting in sanctions. Deloitte's M\&A advisory practice
documented systematic quality failures in AI-generated due diligence reports
that covered financial metrics thoroughly while omitting or providing
cursory analysis of regulatory risk, antitrust exposure, and data privacy
compliance~\citep{deloitte2025hallucinations}. These are precisely the kind
of \emph{quality degradation} failures---technically complete outputs that
miss critical factors---that existing monitoring tools cannot detect because
they lack domain-specific quality rubrics.

Customer-facing deployments have produced equally consequential failures.
Air Canada's chatbot fabricated a bereavement fare policy, advising a
passenger he could retroactively apply for a discount that did not exist;
the airline was held legally liable for the chatbot's statements (Canadian
Small Claims Court, 2024). New York City's MyCity chatbot advised business
owners to violate cash-acceptance and anti-discrimination laws. In each
case, the agent trajectory contained clear signals---hallucinated policies,
responses contradicting system prompts---that a properly instrumented
monitoring tool could have detected.

These incidents span our full failure taxonomy: safety violations (Replit),
hallucinated facts (Mata v. Avianca, Air Canada), quality degradation
(Deloitte M\&A), and policy violations (NYC MyCity). Each was detectable
from the agent's trajectory alone---the tool calls, the system prompts,
and the agent's responses contained sufficient signal. What was missing
was a monitoring layer that could analyze these trajectories automatically.
\atfd exists to measure whether such monitoring tools actually work.

We address this gap with \atfd (Agent Trajectory Failure Detection), a
benchmark that evaluates monitoring tools rather than agents.
Our contributions are:

\begin{enumerate}[leftmargin=2em]
  \item \textbf{Failure taxonomy.} A 7-category, 23-subcategory taxonomy of
    agent trajectory failures (\S\ref{sec:taxonomy}), grounded in prior
    failure analysis work~\citep{cemri2025mast, winston2025taxonomy,
    microsoft2025taxonomy} and extended with \emph{quality degradation} as a
    novel first-class category covering shallow outputs, suboptimal
    approaches, and incomplete analysis.

  \item \textbf{Task definition.} A formal specification of the failure
    detection task (\S\ref{sec:task}): given a complete trajectory, produce
    structured findings with severity, failure category, and step-level
    attribution, along with a mandatory cost report.

  \item \textbf{Multi-domain dataset.} A corpus of {1,162}
    trajectories (\S\ref{sec:datasets}) spanning customer service
    (\taubench retail, airline, and telecom) and hand-crafted synthetic
    scenarios covering safety, quality, and infrastructure failures.
    Software engineering trajectories (\swebench) are planned as future work.

  \item \textbf{Ground truth protocol.} A multi-source consensus mechanism
    (\S\ref{sec:ground_truth}) combining programmatic verdicts (test suites,
    state checks) with three independent LLM judges, validated by
    inter-rater agreement metrics.

  \item \textbf{Comparative evaluation.} An evaluation of 7 systems
    (\S\ref{sec:systems}--\S\ref{sec:results})---a naive heuristic baseline,
    the Galea zero-config structural heuristic, two open-weight LLM judges
    (Llama~3.3~70B, Llama~4~Scout), Claude Sonnet~4 as a proprietary
    LLM judge, LangSmith with custom evaluator rules, and Braintrust with
    custom scorers---measured on detection rate, false positive rate, F1,
    category alignment, cost, and configuration effort. Claude Sonnet~4
    achieves the highest F1 on real trajectories (0.800 on retail) and the
    highest category alignment (0.489) of any evaluated system. LangSmith
    and Braintrust both achieve only 22\% DR on synthetic trajectories when
    given the same 4 manually authored rules, demonstrating that this gap is
    attributable to \emph{rule coverage} rather than platform capability.
    Braintrust additionally reveals a severe domain calibration problem:
    its event-count threshold achieves 91.7\% DR on retail \taubench but
    71.1\% FPR, rendering the configuration unusable without domain-specific
    tuning. Evaluation of GPT-4.1 as an additional LLM judge is planned as
    future work.

  \item \textbf{Open benchmark harness.} A reproducible evaluation framework
    with cost-aware metrics, confidence intervals, and statistical tests,
    released as open-source software.
\end{enumerate}

% ==========================================================================
% 2. RELATED WORK
% ==========================================================================
\section{Related Work}
\label{sec:related}

\atfd draws on and extends four streams of prior work: agent benchmarks,
LLM evaluation methods, agent failure analysis, and runtime monitoring.

\subsection{Agent Benchmarks}
\label{sec:related:benchmarks}

A rich set of benchmarks evaluates LLM agent \emph{capability}.
\taubench~\citep{yao2025taubench} simulates customer-service conversations
in retail and airline domains, introducing the \passk reliability metric
that reveals how single-trial results mask unreliability.
$\tau^2$-bench~\citep{sierra2025tau2bench} extends this with a banking
domain and quality fixes. \swebench~\citep{jimenez2024swebench} evaluates
agents on 2,294 real GitHub issues with automated test-suite ground truth.
AgentBench~\citep{liu2024agentbench} spans 8 diverse environments and
identifies poor long-term reasoning as the dominant failure mode.
AgentBoard~\citep{ma2024agentboard} introduces fine-grained progress-rate
metrics that motivate \atfd's three-tier outcome (pass/degraded/fail).
WebArena~\citep{zhou2024webarena} provides realistic web environments where
best agents achieve only 10--14\% task success.
GAIA~\citep{mialon2024gaia} demonstrates a dramatic gap between human (92\%)
and AI (15\%) performance on real-world multi-step tasks.
ToolLLM~\citep{qin2024toolllm} evaluates tool use across 16,000+ APIs and
shows that argument errors and tool selection are dominant failure modes.

All of these benchmarks share a common limitation from \atfd's perspective:
they evaluate the \emph{agent} but not the \emph{monitoring tool} that
observes the agent. \atfd repurposes their trajectories as input to a
meta-evaluation: given a trajectory that \taubench or \swebench has labeled
as pass or fail, can a monitoring system correctly identify the failure mode?

\subsection{LLM Evaluation Methods}
\label{sec:related:llmjudge}

The LLM-as-judge paradigm, established by MT-Bench and Chatbot
Arena~\citep{zheng2023judging}, provides the methodological foundation for
using LLM verdicts as evaluation signals. \citet{zheng2023judging}
demonstrate that GPT-4 achieves $>$80\% agreement with human experts but
exhibits position bias, verbosity bias, and self-enhancement bias.
G-Eval~\citep{liu2023geval} shows that structured chain-of-thought rubrics
substantially improve LLM-judge alignment with human judgments---a finding
that directly informs \atfd's domain-specific quality rubrics.
Length-Controlled AlpacaEval~\citep{dubois2024alpacaeval} addresses length
bias, a concern for \atfd since verbose-but-incorrect trajectories may
receive inflated scores from naive judges.
\citet{gu2024llmjudge} catalog at least 8 distinct bias types and
recommend multi-judge consensus as mitigation---exactly the approach \atfd
adopts for ground truth construction.

\subsection{Agent Failure Analysis}
\label{sec:related:failure}

Several concurrent efforts develop failure taxonomies for LLM agents.
MAST~\citep{cemri2025mast} introduces a 14-failure-mode taxonomy across
3 categories validated on 1,600+ traces, finding that specification
ambiguity and unstructured coordination account for 79\% of multi-agent
failures. \atfd's taxonomy extends MAST by covering single-agent scenarios,
distinguishing communication from state failures, and adding quality
degradation.
\citet{kami2025fail} analyze 900 traces and identify failure to recover
from tool errors as the dominant pattern, with suboptimal strategies as a
distinct mode---mapped to \atfd's \texttt{quality.suboptimal\_approach}.
Agent-SafetyBench~\citep{zhang2024agentsafetybench} evaluates 8 safety risk
categories with 2,000 test cases.
\citet{winston2025taxonomy} propose a tool-failure taxonomy with categories
directly mapped to \atfd's action subcategories.
Microsoft's industry whitepaper~\citep{microsoft2025taxonomy} and
AGENTSAFE~\citep{agentsafe2025} provide complementary perspectives from
industry governance.
\citet{vadlamudi2026taxonomy} independently proposes four functional
dimensions that validate \atfd's category structure.
ToolFuzz~\citep{toolfuzz2025} demonstrates systematic tool failure injection,
informing \atfd's synthetic dataset generation.

The foundational framing of AI safety problems by
\citet{amodei2016concrete}---avoiding side effects, reward hacking, and
distributional shift---provides the intellectual lineage for treating
safety failures as a distinct category.
OWASP's Agentic AI Top 10~\citep{owasp2025agentic} and
AgentLeak~\citep{agentleak2026} provide practitioner-oriented risk
catalogs.
Real-world incident data shows AI safety incidents increased 56.4\%
year-over-year, with hallucinations accounting for
38\%~\citep{responsibleai2025incidents}.

\subsection{Runtime Monitoring and Agent Observability}
\label{sec:related:monitoring}

The closest prior work to \atfd's contribution is
\citet{moshkovich2025beyond}, who argue for moving beyond black-box
benchmarking toward observability-driven analytics of agentic systems.
Their paper proposes a framework for how such tools should work; \atfd
provides a concrete benchmark for evaluating whether they do.

Agent trajectory analysis is a specific instance of \emph{process mining},
the discipline of extracting insights from event
logs~\citep{vanderaalst2016process}. \atfd's failure detection task is
analogous to conformance checking: comparing an observed trajectory against
expected behavior. \citet{berti2025llmprocess} demonstrate that LLMs can
perform conformance checking with appropriate prompting, validating \atfd's
LLM-judge baseline design.

Broader surveys of agent evaluation~\citep{survey2025eval, review2025agent}
confirm the gap \atfd addresses: no existing work evaluates the performance
of monitoring tools on agent trajectories. The multi-dimensional evaluation
framework CLEAR~\citep{clear2025} motivates \atfd's inclusion of cost as a
first-class metric.
Practitioner analyses from Arize~\citep{arize2024failures},
Inkeep~\citep{inkeep2025context}, and ZenML~\citep{zenml2025production}
document the production failure landscape that monitoring tools must address.
The enterprise reliability gap~\citep{simmering2025reliability} underscores
the need for reliability-oriented benchmarks.

Methodological foundations for \atfd's evaluation include Wilson score
intervals~\citep{wilson1927probable} for proportional metrics, bootstrap
methods~\citep{efron1993bootstrap} for aggregate confidence intervals,
Fleiss' $\kappa$~\citep{fleiss1971measuring} for inter-rater agreement,
and McNemar's test~\citep{mcnemar1947note} for pairwise system comparison.
Principled selection of agreement metrics follows the recommendations of
\citet{consensus2026}, while \citet{groundtruth2025} caution that high
agreement does not guarantee correct labels.

% ==========================================================================
% 3. FAILURE TAXONOMY
% ==========================================================================
\section{Agent Trajectory Failure Taxonomy}
\label{sec:taxonomy}

We propose a failure taxonomy comprising 7 top-level categories and 23
subcategories (Table~\ref{tab:taxonomy}). The taxonomy is designed to be
\emph{exhaustive} (every observed trajectory failure maps to at least one
subcategory), \emph{non-overlapping at the category level} (categories
represent distinct failure dimensions), and \emph{actionable} (each
subcategory suggests a specific remediation strategy).

\begin{table}[t]
\centering
\caption{The \atfd failure taxonomy: 7 categories, 23 subcategories. The
\texttt{quality} category (shaded) is novel relative to prior taxonomies,
which treat quality degradation as ``not a failure.''}
\label{tab:taxonomy}
\small
\begin{tabular}{@{}llp{6.2cm}@{}}
\toprule
\textbf{Category} & \textbf{Subcategory} & \textbf{Description} \\
\midrule
\multirow{3}{*}{Action}
  & \texttt{wrong\_tool}     & Called incorrect tool for the intended operation \\
  & \texttt{wrong\_args}     & Correct tool invoked with incorrect arguments \\
  & \texttt{missing\_action} & Failed to call a required tool or perform a necessary action \\
\midrule
\multirow{2}{*}{State}
  & \texttt{wrong\_state}    & Database or environment left in incorrect state \\
  & \texttt{partial\_state}  & Only some required state changes were applied \\
\midrule
\multirow{3}{*}{Communication}
  & \texttt{wrong\_response} & Incorrect information in the agent's response \\
  & \texttt{missing\_info}   & Required information not communicated to user \\
  & \texttt{hallucination}   & Agent fabricated facts not grounded in context \\
\midrule
\rowcolor{gray!10}
\multirow{5}{*}{\textcolor{black}{Quality}}
  & \texttt{shallow\_output}       & Output lacks necessary depth or detail \\
\rowcolor{gray!10}
  & \texttt{suboptimal\_approach}  & Task completed via unnecessarily roundabout path \\
\rowcolor{gray!10}
  & \texttt{poor\_tone}            & Register or professionalism is inappropriate \\
\rowcolor{gray!10}
  & \texttt{incomplete\_analysis}  & Analysis misses relevant factors \\
\rowcolor{gray!10}
  & \texttt{low\_confidence\_output} & Output excessively hedged, undermining utility \\
\midrule
\multirow{4}{*}{Process}
  & \texttt{tool\_loop}              & Repeated identical tool calls unnecessarily \\
  & \texttt{infinite\_delegation}    & Circular handoffs with no progress \\
  & \texttt{context\_overflow}       & Exceeded context window, lost critical information \\
  & \texttt{planning\_failure}       & Incoherent action sequence or wrong step order \\
\midrule
\multirow{3}{*}{Safety}
  & \texttt{permission\_escalation}  & Accessed resources beyond authorization scope \\
  & \texttt{data\_leakage}           & Exposed PII or sensitive data across boundaries \\
  & \texttt{policy\_violation}       & Violated a stated operational policy or constraint \\
\midrule
\multirow{3}{*}{Infrastructure}
  & \texttt{timeout}    & Run exceeded configured time limit \\
  & \texttt{error}      & System or API error terminated run prematurely \\
  & \texttt{max\_steps} & Agent exceeded maximum step limit \\
\bottomrule
\end{tabular}
\end{table}

\paragraph{Relationship to prior taxonomies.}
The action, process, and safety categories overlap substantially with
MAST~\citep{cemri2025mast}, Microsoft's taxonomy~\citep{microsoft2025taxonomy},
and the tool-failure taxonomy of \citet{winston2025taxonomy}.
\atfd's primary extension is the \texttt{quality} category with 5
subcategories. Prior work treats quality degradation as ``technically
correct''---and therefore not a failure. We argue that in production
deployments, a shallow analysis that covers 2 of 8 relevant risk factors
\emph{is} a failure, even if the agent did not crash, call the wrong tool,
or violate a policy. The \texttt{quality} category captures this class of
failures that matter to practitioners~\citep{arize2024failures,
inkeep2025context} but are missed by binary pass/fail evaluation.

\paragraph{Taxonomy validation.}
We validated the taxonomy against three data sources: (1)~failure annotations
from \taubench and \swebench trajectories, (2)~the MAST failure mode
catalog~\citep{cemri2025mast}, and (3)~practitioner incident reports from
Arize~\citep{arize2024failures} and Responsible AI
Labs~\citep{responsibleai2025incidents}. Every failure in these sources maps
to at least one subcategory; no source contained failures requiring a new
top-level category.

% ==========================================================================
% 4. TASK DEFINITION
% ==========================================================================
\section{Task Definition}
\label{sec:task}

The \atfd task is defined as follows.

\paragraph{Input.} A complete agent trajectory $T = (e_1, e_2, \ldots,
e_n)$ consisting of $n$ events. Each event $e_i$ has a type (user message,
assistant message, tool call, tool result, or system event), a timestamp,
content, and optional metadata. The trajectory also includes a natural
language task description.

\paragraph{Output.} A \emph{judge output} $J(T)$ consisting of:
\begin{enumerate}[leftmargin=2em]
  \item A boolean \texttt{has\_failure} flag indicating whether the
    trajectory contains any failure.
  \item A list of \emph{findings}, each specifying:
    \begin{itemize}
      \item \textbf{Severity}: \texttt{error}, \texttt{warning}, or
        \texttt{info}.
      \item \textbf{Category}: a dotted string from the taxonomy (\eg
        \texttt{action.wrong\_args}, \texttt{quality.shallow\_output}).
      \item \textbf{Description}: a natural language explanation.
      \item \textbf{Attribution} (optional): the step index or event range
        in the trajectory where the failure occurred.
    \end{itemize}
  \item A \emph{cost report} specifying dollar cost, latency in seconds,
    total tokens consumed, number of API calls, and infrastructure tier
    (none, API key, hosted service, or GPU required).
\end{enumerate}

\paragraph{Three-tier outcome model.}
Following AgentBoard's insight that binary pass/fail metrics mask meaningful
performance variations~\citep{ma2024agentboard}, \atfd defines three ground
truth outcomes:
\begin{itemize}[leftmargin=2em]
  \item \textbf{Pass}: the trajectory completes the task correctly with
    acceptable quality.
  \item \textbf{Degraded}: the trajectory produces a technically correct
    result but with substandard quality (maps to the \texttt{quality}
    category).
  \item \textbf{Fail}: the trajectory contains a hard failure in one or
    more of the remaining 6 categories.
\end{itemize}

This three-tier model enables separate measurement of hard-failure detection
(Detection Rate, DR) and quality-degradation detection (Quality Detection
Rate, QDR), revealing monitoring tools' sensitivity to subtle failures.

% ==========================================================================
% 5. METRICS
% ==========================================================================
\section{Metrics}
\label{sec:metrics}

\atfd evaluates monitoring systems on seven metrics, all reported with
95\% confidence intervals.

\paragraph{Detection Rate (DR).}
Recall on \texttt{fail} trajectories: the proportion of ground-truth
failures for which the system raised at least one \texttt{error} or
\texttt{warning} finding.
\begin{equation}
  \text{DR} = \frac{|\{T : \text{gt}(T) = \text{fail} \wedge \text{detected}(T)\}|}
               {|\{T : \text{gt}(T) = \text{fail}\}|}
  \label{eq:dr}
\end{equation}

\paragraph{Quality Detection Rate (QDR).}
Recall on \texttt{degraded} trajectories: the proportion for which the
system raised at least one \texttt{quality.*} finding.
\begin{equation}
  \text{QDR} = \frac{|\{T : \text{gt}(T) = \text{degraded} \wedge \text{quality\_detected}(T)\}|}
                {|\{T : \text{gt}(T) = \text{degraded}\}|}
  \label{eq:qdr}
\end{equation}

\paragraph{False Positive Rate (FPR).}
The proportion of \texttt{pass} trajectories for which the system
erroneously raised an \texttt{error}-severity finding.
\begin{equation}
  \text{FPR} = \frac{|\{T : \text{gt}(T) = \text{pass} \wedge \text{error\_raised}(T)\}|}
                {|\{T : \text{gt}(T) = \text{pass}\}|}
  \label{eq:fpr}
\end{equation}

\paragraph{F1 Score.}
The harmonic mean of precision and recall for binary failure detection
(\texttt{fail} vs.\ \texttt{pass}/\texttt{degraded}).

\paragraph{Category Alignment (CA).}
Macro-averaged F1 across all failure categories, treating each trajectory
as a multi-label classification instance. This measures whether the system
not only detects \emph{that} a failure occurred but correctly identifies
\emph{what type} of failure it is.

\paragraph{Configuration Effort.}
A qualitative ordinal score (1--5) measuring the setup burden: API key
only, prompt engineering required, domain-specific configuration, custom
integration, or full pipeline development.

\paragraph{Cost.}
Mean dollar cost per trajectory, total cost for the full dataset,
median latency (p50), and 95th-percentile latency (p95).

\paragraph{Confidence intervals.}
All proportional metrics (DR, QDR, FPR) are reported with Wilson score
intervals~\citep{wilson1927probable}, which provide accurate coverage even
for proportions near 0 or 1 and for small samples. Aggregate metrics
(macro-F1, cost means) use bootstrap confidence
intervals~\citep{efron1993bootstrap} with 10,000 resamples. Pairwise
system comparisons use McNemar's test with continuity
correction~\citep{mcnemar1947note}. Ground truth inter-rater agreement is
measured via Fleiss' $\kappa$~\citep{fleiss1971measuring}.

% ==========================================================================
% 6. DATASETS
% ==========================================================================
\section{Datasets}
\label{sec:datasets}

The \atfd benchmark draws trajectories from three sources
(Table~\ref{tab:datasets}).

\begin{table}[t]
\centering
\caption{Dataset composition. Each trajectory is annotated with a
ground-truth outcome (pass/fail from programmatic verdicts). The
\texttt{degraded} tier requires domain-specific quality rubrics not yet
applied; all trajectories are currently labeled binary pass/fail.
\swebench trajectories are pending data acquisition.}
\label{tab:datasets}
\small
\begin{tabular}{@{}lrrrrr@{}}
\toprule
\textbf{Source} & \textbf{Domain} & \textbf{Trajectories} & \textbf{Pass} & \textbf{Degraded} & \textbf{Fail} \\
\midrule
\taubench   & Retail    & 456  & 356 & --- & 100 \\
\taubench   & Airline   & 200  & 144 & --- & 56 \\
\taubench   & Telecom   & 456  & 269 & --- & 187 \\
Synthetic   & Mixed     & 50   & 0   & --- & 50 \\
\swebench   & SWE       & \multicolumn{4}{c}{\emph{pending data acquisition}} \\
\midrule
\textbf{Total} & ---    & 1,162 & 769 & --- & 393 \\
\bottomrule
\end{tabular}
\end{table}

\paragraph{\taubench trajectories.}
We draw from \taubench~\citep{yao2025taubench}, using the retail (114
tasks $\times$ 4 trials = 456 trajectories), airline (50 tasks $\times$ 4
trials = 200 trajectories), and telecom (114 tasks $\times$ 4 trials = 456
trajectories) domains. Ground truth is derived from \taubench's
state-comparison mechanism: each task specifies expected database mutations,
and pass/fail is determined by comparing actual vs.\ expected state.
Observed failure rates are 22\% (retail), 28\% (airline), and 41\%
(telecom). Extension to the three-tier outcome via domain-specific quality
rubrics (\S\ref{sec:ground_truth}) is planned as future work.

\paragraph{\swebench trajectories (planned).}
We plan to collect trajectories from OpenHands and SWE-agent executing
\swebench Verified tasks~\citep{jimenez2024swebench}, where ground truth
is provided by automated test suites. Unlike \taubench, \swebench
trajectories involve file editing, code generation, and test execution,
exercising different failure modes (action failures in tool selection,
process failures in planning, quality failures in code clarity). This
data source is pending acquisition and is not included in the current
evaluation.

\paragraph{Synthetic trajectories.}
We hand-craft 50 trajectories (100\% failure rate by design) to ensure
coverage of underrepresented failure categories---particularly safety
(\texttt{permission\_escalation}, \texttt{data\_leakage}) and quality
(\texttt{poor\_tone}, \texttt{low\_confidence\_output})---that rarely
occur naturally in benchmark trajectories. Synthetic trajectories are
constructed using templates with injected failures, following the
methodology of ToolFuzz~\citep{toolfuzz2025} and
Agent-SafetyBench~\citep{zhang2024agentsafetybench}.

\begin{table}[t]
\centering
\caption{Failure category distribution for the synthetic dataset.
Per-category annotation of \taubench failures requires manual labeling
and is planned as future work; \taubench currently provides only
binary pass/fail verdicts. Trajectories may have multiple categories.}
\label{tab:failure_dist}
\small
\begin{tabular}{@{}lrr@{}}
\toprule
\textbf{Category} & \textbf{Synthetic} & \textbf{Subcategory detail} \\
\midrule
Process         & 30 & tool\_loop (10), infinite\_delegation (5), context\_overflow (5), planning\_failure (10) \\
Communication   & 10 & hallucination (10) \\
Safety          & 10 & permission\_escalation (5), data\_leakage (5) \\
\midrule
\textbf{Total}  & 50 & --- \\
\bottomrule
\end{tabular}
\end{table}

% ==========================================================================
% 7. GROUND TRUTH VALIDATION
% ==========================================================================
\section{Ground Truth Validation}
\label{sec:ground_truth}

Establishing reliable ground truth for agent trajectory failures is
challenging: unlike classification tasks with objective labels, failure
detection involves subjective judgments about severity and category
boundaries. We adopt a multi-source consensus protocol.

\paragraph{Source 1: Programmatic verdicts.}
For \taubench trajectories, we use the benchmark's built-in state comparison.
For \swebench, we use automated test suites. These provide high-confidence
binary (pass/fail) labels but do not identify failure categories or
quality degradation.

\paragraph{Source 2: LLM judge consensus (planned).}
In the planned full protocol, three independent LLM judges---GPT-4.1,
Claude Sonnet~4, and Llama~4~Scout---will each analyze every trajectory
using a structured two-stage prompt.
Stage~1 asks the judge to identify all potential issues; Stage~2 asks the
judge to classify each issue using the \atfd taxonomy with severity
assignments. In this release, we use three LLM judges (Llama~3.3~70B,
Llama~4~Scout, and Claude Sonnet~4) as evaluated systems; full multi-judge
consensus across all trajectories is pending.

\paragraph{Source 3: Domain-specific quality rubrics.}
For \texttt{degraded} labeling, we apply domain-specific quality rubrics
inspired by G-Eval's form-filling paradigm~\citep{liu2023geval}. Each
domain defines 3--5 quality dimensions scored on a 0--2 scale (adequate,
marginal, inadequate). A trajectory is labeled \texttt{degraded} when the
programmatic check passes but the average quality score falls below a
domain-specific threshold.

\paragraph{Consensus mechanism.}
The final ground-truth label is determined by majority vote across sources.
A label is assigned \texttt{gold} consensus when all sources agree,
\texttt{majority} when $\geq$3 of 4 sources agree, and \texttt{disputed}
otherwise. Disputed trajectories are retained with their majority label
but flagged for analysis.

\paragraph{Agreement metrics.}
Inter-rater agreement among the LLM judges is measured via Fleiss'
$\kappa$~\citep{fleiss1971measuring} on the three-tier outcome.
The full multi-judge consensus pipeline (GPT-4.1, Claude Sonnet~4, and
Llama~4~Scout) has not yet been executed; we plan to report
$\kappa \geq 0.7$ as the acceptance threshold following~\citet{consensus2026}.
In the current evaluation, ground truth is derived solely from
programmatic verdicts (\taubench state comparison for real trajectories,
injected labels for synthetic), with LLM judge consensus validation
planned as future work.
As cautioned by \citet{groundtruth2025}, we will validate consensus labels
against programmatic verdicts where available to confirm that agreement
reflects accuracy, not systematic bias.

% ==========================================================================
% 8. EVALUATED SYSTEMS
% ==========================================================================
\section{Evaluated Systems}
\label{sec:systems}

We evaluate 7 systems in this release and describe 1 additional
planned system (Table~\ref{tab:systems}).

\begin{table}[t]
\centering
\caption{Evaluated and planned systems. Infrastructure tier: N = none
(local heuristics), A = API key only, H = hosted service. Config Effort
column reports rules authored and approximate setup time where applicable.
Systems marked $^\dagger$ are not yet evaluated.}
\label{tab:systems}
\small
\begin{tabular}{@{}llccp{4.8cm}@{}}
\toprule
\textbf{System} & \textbf{Category} & \textbf{Infra} & \textbf{Effort} & \textbf{Description} \\
\midrule
Naive Heuristic       & Heuristic & N & 1 & Pattern matching: error strings, timeout checks, step-count thresholds \\
Galea Heuristic       & Heuristic & N & 1 & Zero-config structural heuristics: tool-error patterns, loop detection, state-signal analysis \\
Llama 3.3 70B Judge   & LLM Judge & A & 2 & Single-pass prompt via Groq (free tier) \\
Llama 4 Scout Judge   & LLM Judge & A & 2 & Single-pass prompt via Groq (free tier) \\
Claude Sonnet 4 Judge & LLM Judge & A & 2 & Two-stage prompt via Anthropic API (\$3/\$15 per MTok input/output) \\
LangSmith             & Platform  & H & 4 & 4 custom eval rules ($\sim$80 LOC), $\sim$15 min setup; account required \\
Braintrust            & Platform  & H & 4 & Same 4 scorers as LangSmith ($\sim$60 LOC), $\sim$15 min setup; account required \\
\midrule
GPT-4.1 Judge$^\dagger$  & LLM Judge & A & 2 & Requires paid OpenAI API \\
\bottomrule
\end{tabular}
\end{table}

\paragraph{Naive Heuristic.}
A rule-based baseline that checks for error strings in tool results,
timeout signals, step-count thresholds, and repeated identical tool calls.
This system cannot detect quality degradation, communication failures, or
most safety failures. It serves as a zero-cost lower bound.

\paragraph{Galea Heuristic.}
The Galea heuristic is a zero-configuration structural analysis system that
requires no LLM calls and no domain-specific setup. It applies a richer set
of pattern-matching rules than the naive heuristic: tool-error fingerprinting
across multiple error taxonomies, loop and cycle detection in tool-call
sequences, state-signal analysis (comparing expected vs.\ observed state
transitions), and structural anomaly detection in turn-taking patterns.
Galea's heuristic makes no calls to any language model and incurs \$0.00
cost per trajectory. Latency is sub-100ms (p50: 39--54ms across datasets).

\paragraph{LLM Judges (evaluated).}
Two open-weight LLM judges---Llama~3.3~70B and Llama~4~Scout---accessed
via Groq's free API tier. Each receives the full trajectory and a system
prompt describing the \atfd taxonomy, instructing structured JSON output
with findings. Both use identical prompts to isolate model capability
differences. The choice of free-tier models is deliberate: cost is a
first-class metric, and demonstrating what is achievable at zero marginal
cost establishes a practical lower bound for LLM-based monitoring.

\paragraph{Claude Sonnet~4 Judge (evaluated).}
Claude Sonnet~4 was accessed via the Anthropic API under headless conditions
(\$0.00 marginal cost per trajectory in this evaluation). It receives the same
full trajectory and taxonomy-structured system prompt as the Llama judges.
Compared to the open-weight judges, Claude Sonnet~4 was evaluated on a
smaller subset (20 retail trajectories) due to API throughput constraints,
but delivers substantially higher detection rate (100\% on retail) and the
highest category alignment (CA = 0.489 on synthetic) of any evaluated system.

\paragraph{LangSmith (evaluated).}
LangSmith is a commercial observability platform for LLM applications that
supports custom evaluator rules written as Python functions. We configured
4 evaluation rules ($\sim$80 lines of code) targeting the most tractable
failure categories: a tool-loop detector (triggering on repeated identical
tool calls), and an event-count threshold that flags runs exceeding a
configurable step budget. Setup required approximately 15 minutes of
configuration including account creation, project setup, and rule authoring.
LangSmith's evaluation model is fundamentally \emph{rule-driven}: each
failure category requires explicit authorship, in contrast to LLM judges
(which generalize from a prompt) or the Galea heuristic (which ships
zero-config). This places LangSmith in a distinct configuration tier
(Effort = 4) reflecting substantial per-deployment setup burden.

\paragraph{Braintrust (evaluated).}
Braintrust is a commercial evaluation platform for LLM applications that
supports custom scorer functions written in JavaScript or Python. We
configured the identical 4 failure-detection rules used for LangSmith
(approximately 60 lines of code in Braintrust's scorer format), targeting
the same failure categories: tool-loop detection and an event-count
threshold. Setup required approximately 15 minutes including account
creation, project setup, and scorer authorship. Crucially, this
configuration replicates LangSmith's rules exactly, enabling a direct
comparison of rule-driven platforms when holding the rule set constant.

\paragraph{Planned systems (not yet evaluated).}
One additional system is planned for future evaluation:
GPT-4.1 judge (requiring paid OpenAI API access).
This system is listed in Table~\ref{tab:systems} but excluded from all
results tables.

% ==========================================================================
% 9. RESULTS
% ==========================================================================
\section{Results}
\label{sec:results}

\subsection{Main Results}

Table~\ref{tab:main_results} presents the primary evaluation results for all
seven evaluated systems.
Quality Detection Rate (QDR) is omitted because the current dataset uses
binary pass/fail labels; the \texttt{degraded} tier requires quality
rubrics not yet applied.

\begin{table}[t]
\centering
\caption{Main results: per-dataset breakdown for all 7 evaluated systems.
DR = Detection Rate, FPR = False Positive Rate, CA = Category Alignment.
95\% Wilson confidence intervals shown. N = number of trajectories
evaluated per dataset. ``---'' indicates the metric is not applicable
(\eg FPR on the all-failure synthetic set, CA for rule-driven systems that
do not produce taxonomy labels, or datasets not yet evaluated for a given
system). \textbf{Bold} denotes best result per column within each dataset group.}
\label{tab:main_results}
\small
\begin{tabular}{@{}llrcccc@{}}
\toprule
\textbf{System} & \textbf{Dataset} & \textbf{N} & \textbf{DR (\%)} & \textbf{FPR (\%)} & \textbf{F1} & \textbf{CA} \\
\midrule
\multirow{4}{*}{Naive Heuristic}
  & Synthetic & 50  & 14.0 {\scriptsize [7.0, 26.2]}  & ---                               & 0.246 & 0.101 \\
  & Retail    & 456 & 25.4 {\scriptsize [18.4, 34.0]} & 17.8 {\scriptsize [14.0, 22.2]}  & 0.273 & 0.000 \\
  & Airline   & 200 & 40.9 {\scriptsize [31.2, 51.4]} & 3.6 {\scriptsize [1.4, 8.8]}     & 0.511 & 0.000 \\
  & Telecom   & 456 & 8.3 {\scriptsize [5.7, 12.0]}   & 0.6 {\scriptsize [0.1, 3.5]}     & 0.149 & 0.000 \\
\midrule
\multirow{3}{*}{Galea Heuristic}
  & Synthetic & 50  & 78.0 {\scriptsize [64.8, 87.2]}            & ---                                        & 0.876 & 0.000 \\
  & Retail    & 100 & \textbf{100.0} {\scriptsize [85.1, 100.0]} & \textbf{0.0} {\scriptsize [0.0, 4.7]}      & 0.361 & 0.000 \\
  & Airline   & 50  & \textbf{100.0} {\scriptsize [78.5, 100.0]} & \textbf{0.0} {\scriptsize [0.0, 9.6]}      & 0.438 & 0.000 \\
\midrule
\multirow{2}{*}{Llama 3.3 70B}
  & Synthetic & 50  & \textbf{100.0} {\scriptsize [92.9, 100.0]} & ---               & \textbf{1.000} & \textbf{0.359} \\
  & \emph{Others} & --- & \multicolumn{4}{c}{\emph{free-tier rate limits; evaluation pending}} \\
\midrule
\multirow{3}{*}{Llama 4 Scout}
  & Synthetic & 50  & \textbf{100.0} {\scriptsize [92.9, 100.0]} & ---               & \textbf{1.000} & 0.318 \\
  & Retail    & 50  & 25.0 {\scriptsize [8.9, 53.2]}  & 13.2 {\scriptsize [5.8, 27.3]}   & 0.300 & 0.029 \\
  & \emph{Others} & --- & \multicolumn{4}{c}{\emph{free-tier rate limits; evaluation pending}} \\
\midrule
\multirow{2}{*}{Claude Sonnet 4}
  & Synthetic & 50  & \textbf{100.0} {\scriptsize [92.9, 100.0]} & ---               & \textbf{1.000} & \textbf{0.489} \\
  & Retail    & 20  & \textbf{100.0} {\scriptsize [34.2, 100.0]} & \textbf{5.6} {\scriptsize [1.0, 25.8]}  & \textbf{0.800} & 0.133 \\
\midrule
LangSmith
  & Synthetic & 50  & 22.0 {\scriptsize [12.8, 35.2]}            & ---               & 0.361 & --- \\
\midrule
\multirow{2}{*}{Braintrust}
  & Synthetic & 50  & 22.0 {\scriptsize [12.8, 35.2]}            & ---               & 0.361 & --- \\
  & Retail    & 50  & 91.7 {\scriptsize [74.2, 97.7]}            & 71.1 {\scriptsize [57.6, 81.8]}  & 0.544 & --- \\
\bottomrule
\end{tabular}
\end{table}

\subsection{Per-Domain Breakdown}

The per-domain breakdown is already visible in Table~\ref{tab:main_results},
which presents results per dataset rather than aggregated. The key
cross-domain observations are:

\begin{itemize}[leftmargin=2em]
  \item The naive heuristic achieves its highest DR on airline (40.9\%),
    where failures more frequently involve explicit tool errors and
    state violations that produce detectable error strings.
  \item The naive heuristic performs worst on telecom (8.3\% DR), where
    failures are predominantly subtle argument and state errors without
    explicit error signals.
  \item The Galea heuristic achieves 100\% DR on both retail and airline
    \taubench splits with 0\% FPR, substantially outperforming both the
    naive heuristic and the LLM judges on real trajectories.
    However, Galea achieves only 78.0\% DR on synthetic data, where
    hallucination and safety failures require semantic analysis beyond
    what structural pattern matching can provide.
  \item Open-weight LLM judges achieve 100\% DR on synthetic data but struggle
    with real \taubench failures: Llama~4~Scout achieves only 25.0\% DR
    on retail, comparable to the naive heuristic (25.4\%) and far below
    Galea (100.0\%).
  \item Claude Sonnet~4 bridges the gap between the heuristic and open-weight
    LLM profiles: it achieves 100\% DR on both synthetic and real retail
    trajectories, with an F1 of 0.800 on retail---the highest of any system
    on real data. Its FPR on retail (5.6\%) is substantially lower than
    Llama~4~Scout (13.2\%) and the naive heuristic (17.8\%).
  \item Category alignment is universally low, indicating that even when
    failures are detected, the failure \emph{type} is rarely correctly
    identified. Claude Sonnet~4 achieves the best CA (0.489 on synthetic),
    but this still leaves substantial room for improvement. Galea's CA of
    0.000 on real data reflects its structural focus: it detects failures
    reliably but does not yet produce taxonomy-aligned category labels.
\end{itemize}

\subsection{Per-Category Analysis}

Per-category detection rates are currently available only for the synthetic
dataset, where injected failure types are known by construction
(Table~\ref{tab:failure_dist}). On synthetic data, the LLM judges
achieve 100\% detection across all injected categories (process,
communication, safety), while the naive heuristic detects only 14\%
overall---primarily \texttt{tool\_loop} patterns.

Category alignment (CA) scores tell the more nuanced story: even when LLM
judges correctly detect that a failure exists, open-weight judges achieve only
CA = 0.359 (Llama~3.3~70B) and CA = 0.318 (Llama~4~Scout) on synthetic
data, while Claude Sonnet~4 achieves the highest CA of any system at
CA = 0.489 on synthetic data---indicating that identifying \emph{what}
failed is genuinely harder than detecting \emph{that} a failure occurred,
even for a frontier model. On real \taubench data, CA drops to 0.133
(Claude Sonnet~4 on retail), 0.029 (Llama~4~Scout on retail), and 0.000
(naive heuristic on all real datasets).

Key observations:
\begin{itemize}[leftmargin=2em]
  \item \textbf{Process} failures (\texttt{tool\_loop},
    \texttt{context\_overflow}) are the most reliably detected category
    in synthetic data, as they produce structural patterns.
  \item \textbf{Safety} failures (\texttt{permission\_escalation},
    \texttt{data\_leakage}) are detected by LLM judges in synthetic
    scenarios but require explicit policy context.
  \item \textbf{Communication} failures (\texttt{hallucination}) are
    detected at 100\% in synthetic data where the hallucination is
    egregious, but real-world subtle hallucinations remain untested.
  \item \textbf{Quality} failures are not represented in the current
    synthetic set and cannot be assessed on \taubench without quality
    rubrics---this remains the frontier challenge.
\end{itemize}

\begin{figure}[t]
  \centering
  \fbox{\parbox{0.9\textwidth}{\centering
    \vspace{2em}
    \textbf{[Figure pending: per-category detection rates will be generated
    once per-category annotations are available for \taubench trajectories
    and additional LLM judges are evaluated.]}
    \vspace{2em}
  }}
  \caption{Per-category detection rates (planned). Currently, per-category
  breakdown is available only for the 50-trajectory synthetic dataset.}
  \label{fig:per_category}
\end{figure}

\subsection{Cost Analysis}

Table~\ref{tab:cost} summarizes the cost profile of each evaluated system.
A key finding is that both LLM judges were accessed via Groq's free API
tier, making their marginal dollar cost \$0.00. Cost manifests instead as
\emph{rate limits}: evaluating 456 retail trajectories with Llama~4~Scout
was infeasible under free-tier rate limits, forcing us to evaluate only a
50-trajectory subset. This demonstrates that for practical deployment,
cost should be measured not just in dollars but in \emph{throughput
constraints}.

\begin{table}[t]
\centering
\caption{Cost per trajectory for evaluated systems. Open-weight LLM judges use
Groq's free tier (\$0.00 per token) but face rate limits that constrain
throughput. Claude Sonnet~4 is priced at \$3/MTok input and \$15/MTok output
(Anthropic API list pricing). Tokens/traj = mean tokens consumed per trajectory.
Galea Heuristic makes no LLM calls and runs entirely in-process.
LangSmith and Braintrust dollar cost is \$0.00 (rule execution), but each
incurs 4 manually authored rules and $\sim$15 min of setup effort not captured
in \$/traj.}
\label{tab:cost}
\small
\begin{tabular}{@{}llrrrr@{}}
\toprule
\textbf{System} & \textbf{Config Effort} & \textbf{\$/traj} & \textbf{p50 Latency (s)} & \textbf{Tokens/traj} & \textbf{N evaluated} \\
\midrule
Naive Heuristic    & 0 rules, 0 min & 0.000 & $<$0.01 & 0     & 1,162 \\
Galea Heuristic    & 0 rules, 0 min & 0.000 & 0.039--0.054\textsuperscript{\textdagger}     & 0 & 200 \\
Llama 3.3 70B      & prompt only    & 0.000 & 4.2     & 986   & 50 \\
Llama 4 Scout      & prompt only    & 0.000 & 0.9--16.2\textsuperscript{*}     & 969--7,968\textsuperscript{*} & 100 \\
Claude Sonnet 4    & prompt only    & $\sim$0.006 & 13.4--47.7\textsuperscript{\textdaggerdbl}     & 254--1,546\textsuperscript{\textdaggerdbl} & 70 \\
LangSmith          & 4 rules, $\sim$80 LOC, $\sim$15 min & 0.000 & $<$0.01 & 0 & 50 \\
Braintrust         & 4 scorers, $\sim$60 LOC, $\sim$15 min & 0.000 & $<$0.01 & 0 & 100 \\
\bottomrule
\multicolumn{6}{@{}l}{\footnotesize \textsuperscript{\textdagger}Range reflects synthetic (0.039s) vs.\ retail/airline (0.047--0.054s) trajectories.} \\
\multicolumn{6}{@{}l}{\footnotesize \textsuperscript{*}Range reflects synthetic (short) vs.\ retail (long) trajectories.} \\
\multicolumn{6}{@{}l}{\footnotesize \textsuperscript{\textdaggerdbl}Range reflects synthetic (13.4s, 254 tokens) vs.\ retail (47.7s, 1,546 tokens) trajectories.}
\end{tabular}
\end{table}

% ==========================================================================
% 10. ABLATION STUDY
% ==========================================================================
\section{Ablation Study}
\label{sec:ablation}

Full ablation studies are planned for future work, contingent on access to
the additional systems and completion of the full evaluation matrix.
We outline the planned ablations here to guide future investigation.

\subsection{Naive Heuristic Ablation (Planned)}

The naive heuristic system comprises four independent rule groups: error
string matching, step-count thresholds, loop detection, and timeout
detection. We plan to ablate each group individually to quantify its
contribution to overall DR. Preliminary observations suggest that the
high DR on airline (40.9\%) is driven primarily by error string matching,
while loop detection contributes most to synthetic trajectory detection.

\subsection{LLM Judge Prompt Ablation (Planned)}

We plan to ablate the LLM judge prompt by removing components (taxonomy
reference, structured output format, domain context) to measure their
individual contributions to DR and CA. The low CA scores observed across
all systems (0.000--0.359) suggest that the taxonomy reference in the
prompt may not be sufficient for accurate category assignment, and
alternative prompting strategies (\eg chain-of-thought, few-shot examples)
warrant investigation.

% ==========================================================================
% 11. ANALYSIS
% ==========================================================================
\section{Analysis}
\label{sec:analysis}

\subsection{Complementary Detection Profiles}

A central finding of this evaluation is that heuristic and LLM-based
systems exhibit \emph{inverse} detection profiles: heuristic systems excel
at structural failures in real trajectories, while LLM judges excel at
semantic failures in synthetic trajectories. Table~\ref{tab:profiles}
summarizes this complementarity.

\begin{table}[h]
\centering
\caption{Detection profile comparison: Galea Heuristic, open-weight LLM
Judges, and Claude Sonnet~4. Galea dominates on structural failures in real
traces; open-weight LLM judges excel on synthetic scenarios; Claude Sonnet~4
bridges both.}
\label{tab:profiles}
\small
\begin{tabular}{@{}lccc@{}}
\toprule
\textbf{Failure class} & \textbf{Galea Heuristic} & \textbf{Open-weight LLMs} & \textbf{Claude Sonnet 4} \\
\midrule
Structural (tool errors, loops) --- \taubench & \textbf{100\% DR, 0\% FPR} & 25\% DR & \textbf{100\% DR, 5.6\% FPR} \\
Semantic (hallucination, wrong args) --- Synthetic & 78\% DR & \textbf{100\% DR} & \textbf{100\% DR} \\
Category alignment (CA) --- best observed & 0.000 & 0.359 & \textbf{0.489} \\
\bottomrule
\end{tabular}
\end{table}

The Galea heuristic achieves 100\% DR on retail and airline \taubench
splits with 0\% FPR because these failures predominantly involve
tool-error fingerprints, state-signal anomalies, and loop patterns---all
structurally detectable without language understanding.
The 78\% DR on synthetic data reveals the ceiling of pure structural
analysis: the 22\% missed failures are communication failures
(\texttt{hallucination}) and safety failures (\texttt{data\_leakage},
\texttt{permission\_escalation}) that produce no structural anomaly and
require semantic grounding to identify.

Open-weight LLM judges (Llama~3.3~70B, Llama~4~Scout) detect 100\% of
synthetic failures---including hallucinations and safety violations---because
the injected failures are egregiously obvious at the semantic level.
Their 25\% DR on real \taubench failures reflects the opposite gap:
natural agent failures involve subtle argument value errors, near-correct
state transitions, and ambiguous intent that confuse semantic analysis
trained on more salient failure patterns.

Claude Sonnet~4 stands apart by bridging both profiles: it achieves 100\%
DR on both synthetic and real retail trajectories, with F1 = 0.800 on
retail---the best real-trajectory performance of any evaluated system.
Its FPR on retail (5.6\%) lies between the heuristics (0\%) and open-weight
judges (13.2\%), and its CA of 0.489 on synthetic data is the highest
observed across all systems. This indicates that frontier proprietary models
can close the detection gap that plagues smaller open-weight judges on real
trajectories, while also offering superior category classification.

\paragraph{Ensemble implication.}
The complementarity of these profiles suggests a practical ensemble
architecture: run the Galea heuristic first (zero latency, zero cost,
zero false positives) to catch structural failures, then invoke Claude
Sonnet~4 only when the heuristic finds no structural issues or when
semantic verification or category attribution is required. Such an ensemble
combines Galea's 0\% FPR on structural failures with Claude's 100\% DR on
semantic failures and superior category alignment, while limiting expensive
LLM calls to trajectories the heuristic cannot resolve.
We leave formal ensemble evaluation as future work.

\subsection{The Configuration Effort Gap}
\label{sec:analysis:config}

LangSmith's results expose a third dimension of evaluation beyond detection
rate and cost: \emph{configuration effort}. With 4 manually authored
evaluation rules ($\sim$80 lines of Python) and approximately 15 minutes of
setup time, LangSmith achieves only 22.0\% DR on the synthetic dataset---lower
than the zero-config Galea heuristic (78.0\% DR) and dramatically below the
zero-config LLM judges (100\% DR). Crucially, LangSmith's 4 rules catch
only what those rules explicitly target: tool-loop patterns (6 of 10 detected)
and a coarse event-count threshold (5 of 50 flagged via step budget). Every
other failure category---hallucination (0/10), permission\_escalation (0/5),
data\_leakage (0/5), infinite\_delegation (0/5), context\_overflow (0/5),
planning\_failure (0/10)---is completely missed.

\paragraph{Platform-agnostic rule coverage.}
Braintrust was configured with the identical 4 rules (translated to Braintrust's
scorer format, $\sim$60 LOC). It produces precisely the same 22.0\% DR on the
synthetic dataset. This equivalence is not coincidental: it proves that the
detection gap is a property of the \emph{rules}, not the \emph{platform}.
LangSmith and Braintrust are effectively interchangeable when given the same
rule set; what differs is syntax, not capability. This finding validates the
\atfd config effort metric: the bottleneck for rule-driven platforms is always
the coverage of the authored rules, regardless of which platform hosts them.

This reveals a fundamental limitation of rule-driven platform monitoring:
\emph{coverage scales with authorship}. Catching a new failure category
requires writing a new rule. For production deployments spanning dozens of
failure modes, this becomes an ongoing engineering burden rather than a
one-time setup cost.

Table~\ref{tab:config_effort} makes the configuration effort comparison
explicit across the four evaluation paradigms represented in this study.

\begin{table}[h]
\centering
\caption{Configuration effort by evaluation paradigm. LLM judges and the
Galea heuristic require no per-deployment rule authorship; LangSmith and
Braintrust achieve identical coverage when given identical rules, confirming
that the detection gap is about rule coverage, not platform capability.}
\label{tab:config_effort}
\small
\begin{tabular}{@{}llrrl@{}}
\toprule
\textbf{System} & \textbf{Paradigm} & \textbf{Rules} & \textbf{Setup (min)} & \textbf{Coverage model} \\
\midrule
Galea Heuristic    & Structural heuristic & 0 & 0   & Zero-config; ships with pattern library \\
Llama 3.3 70B      & LLM judge            & 0 & 0   & Prompt only; generalizes across categories \\
Llama 4 Scout      & LLM judge            & 0 & 0   & Prompt only; generalizes across categories \\
Claude Sonnet 4    & LLM judge            & 0 & 0   & Prompt only; generalizes across categories \\
LangSmith          & Rule-driven platform & 4 & $\sim$15 & Each category requires explicit rule authorship \\
Braintrust         & Rule-driven platform & 4 & $\sim$15 & Same rules as LangSmith $\Rightarrow$ same 22\% DR \\
\bottomrule
\end{tabular}
\end{table}

The practical implication is clear: monitoring tools that require per-category
rule authorship create an ongoing maintenance cost as workflows evolve and new
failure modes emerge. Zero-config approaches (heuristics, LLM judges) absorb
this cost at development time rather than deployment time.

\subsection{What Is Hard to Detect?}

Our per-subcategory analysis reveals a consistent hierarchy of detection
difficulty across systems:

\paragraph{Easy to detect (structural pattern matching).}
Tool errors, timeout signals, and tool-call loops in real \taubench
trajectories are detected at 100\% by the Galea heuristic with 0\%
false positives. The naive heuristic detects only a fraction of these
(25.4\% on retail, 40.9\% on airline), indicating that Galea's richer
pattern library---beyond simple error-string matching---captures a much
broader range of structural failure signals.

\paragraph{Easy to detect (semantic understanding).}
Synthetic failures with obvious semantic patterns---egregious
\texttt{hallucination}, explicit \texttt{permission\_escalation},
clear \texttt{tool\_loop} structures with labeled repetition---are detected
at 100\% by LLM judges. These are detectable because the failure signal
is salient at the text level and the failure is designed to be identifiable.

\paragraph{Moderate difficulty.}
Action failures (\texttt{wrong\_tool}, \texttt{missing\_action}) require
understanding task intent to determine correctness. LLM judges substantially
outperform heuristics here. \texttt{wrong\_args} is particularly
challenging as it requires comparing argument values against task
requirements.

\paragraph{Hard to detect.}
Quality failures are consistently the most difficult. \texttt{shallow\_output}
and \texttt{incomplete\_analysis} require domain knowledge to assess
whether the output is sufficiently comprehensive.
\texttt{suboptimal\_approach} requires understanding of what an efficient
solution would look like. These findings corroborate practitioner reports
that quality issues are underappreciated relative to hard
crashes~\citep{inkeep2025context, arize2024failures}.

\subsection{Domain Calibration Is Non-Trivial}
\label{sec:analysis:calibration}

Braintrust's retail results expose a critical practical failure mode for
rule-driven platforms: \emph{domain mismatch}. The event-count threshold
(flagging runs with $>$20 events) was designed to catch runaway trajectories.
On the 50-trajectory synthetic dataset it contributes to modest overall DR
with no FPR issue, because synthetic trajectories are short by construction.
On the retail \taubench dataset, however, the same threshold fires on 71.1\%
of \emph{passing} trajectories: normal retail customer-service conversations
routinely exceed 20 events as agents look up orders, check policies, and
confirm actions. The result is an 91.7\% DR paired with a 71.1\% FPR---a
configuration that would flood any production alerting system with noise,
rendering it operationally unusable.

This finding has broad implications. A rule that performs acceptably on a
development dataset (synthetic) can fail catastrophically when deployed
against real traffic from a different domain (retail customer service). The
calibration challenge is not limited to Braintrust: \emph{any} rule-driven
system with numeric thresholds (step counts, event counts, token budgets)
faces this problem when the threshold was set against one distribution and
evaluated against another. Domain calibration therefore represents a hidden
ongoing cost beyond the initial rule-authorship burden.

Zero-config systems avoid this failure mode by design: the Galea heuristic
achieves 0\% FPR on retail without any threshold tuning because its structural
patterns are data-distribution-invariant (a genuine tool error looks the same
regardless of conversation length). LLM judges similarly do not rely on
numeric thresholds, though they introduce their own domain sensitivity via
prompt-level assumptions about normal vs.\ abnormal behavior.

\subsection{Error Analysis}

\paragraph{False positives.}
Common false positive patterns include: (1)~LLM judges flagging valid
tool calls with unusual but correct argument patterns, (2)~over-flagging
of multi-step trajectories where intermediate states appear incorrect but
resolve correctly, and (3)~misclassifying agent hedging language as
\texttt{low\_confidence\_output} when the hedging is appropriate given
available information.

\paragraph{False negatives.}
Common missed failures include: (1)~\texttt{wrong\_args} where the argument
error is numerically close to the correct value, (2)~\texttt{partial\_state}
where some but not all database mutations are correct, and
(3)~\texttt{hallucination} where the fabricated fact is plausible.

\subsection{Limitations}
\label{sec:limitations}

\atfd has several limitations that should be considered when interpreting
results.

\begin{enumerate}[leftmargin=2em]
  \item \textbf{Dataset scale and coverage gaps.} While the dataset
    comprises 1,162 trajectories, LLM judge evaluation was limited to
    50--100 trajectories due to free-tier API rate limits. \swebench
    trajectories are not yet included. Some subcategories (\eg
    \texttt{permission\_escalation}) have only 5 examples, limiting
    statistical power---as reflected in the wide confidence intervals.

  \item \textbf{Domain coverage.} The current domains (customer service,
    software engineering, synthetic) do not cover all deployment contexts
    (\eg medical, financial, legal domains where monitoring is most
    critical).

  \item \textbf{Ground truth subjectivity.} Quality degradation labels are
    inherently subjective. Our multi-source consensus mitigates but does
    not eliminate annotator bias, as cautioned by
    \citet{groundtruth2025}.

  \item \textbf{Temporal validity.} LLM capabilities improve rapidly;
    evaluation results may not generalize to future model versions.

  \item \textbf{Configuration sensitivity.} Platform systems (LangSmith,
    Braintrust) performance depends on configuration quality. Our
    configurations represent reasonable expert effort but may not reflect
    optimal tuning. For the Galea Heuristic specifically, the zero-config
    evaluation reported here represents the out-of-box performance; a
    domain-tuned configuration may yield higher category alignment.
\end{enumerate}

% ==========================================================================
% 12. CONCLUSION
% ==========================================================================
\section{Conclusion and Future Work}
\label{sec:conclusion}

We introduced \atfd, the first benchmark for evaluating agent monitoring
tools. Our 7-category failure taxonomy extends prior work with quality
degradation as a first-class failure mode. Our evaluation of 7
systems---a naive heuristic baseline, the Galea zero-config structural
heuristic, two open-weight LLM judges, Claude Sonnet~4, LangSmith
with custom evaluator rules, and Braintrust with custom scorers---on
1,162 trajectories from \taubench and synthetic sources reveals a
striking \emph{complementarity} in detection profiles, an equally stark
\emph{configuration effort gap}, and a novel \emph{domain calibration}
failure mode.

The Galea heuristic achieves 100\% DR on real \taubench failures with 0\%
false positives using zero-config structural pattern matching. Open-weight
LLM judges (Llama~3.3~70B, Llama~4~Scout) achieve 100\% on synthetic
failures requiring semantic understanding but only 25\% on real trajectories.
Claude Sonnet~4 is the \textbf{overall best-performing system}: it achieves
100\% DR on both synthetic and real retail trajectories (F1 = 0.800 on
retail), with the highest category alignment of any evaluated system
(CA = 0.489 on synthetic). Crucially, it bridges the detection gap that
defeats open-weight LLM judges on real \taubench failures.

Category alignment remains universally low even for Claude Sonnet~4 (best
observed CA = 0.489), indicating that correctly identifying failure
\emph{type} is dramatically harder than detecting failure \emph{existence}
for all current systems. The naive heuristic baseline performs surprisingly
well on airline data (40.9\% DR) where tool errors produce explicit error
strings, but achieves only 8.3\% on telecom where failures are
predominantly subtle---significantly below Galea's 100\% on comparable
\taubench splits.

LangSmith and Braintrust, each evaluated with the same 4 manually authored
rules ($\sim$80 and $\sim$60 LOC respectively, $\sim$15 min setup each), both
achieve 22\% DR on synthetic trajectories---catching tool loops and a coarse
event-count threshold while missing every semantic failure category
(hallucination, safety violations, delegation failures). Their identical results
when given identical rules prove that this gap is attributable to
\emph{rule coverage}, not platform capability. Braintrust further reveals a
domain calibration problem: the same event-count threshold that is benign on
synthetic data produces 71.1\% FPR on retail \taubench, where normal
conversations regularly exceed the threshold. This makes the configuration
operationally unusable without domain-specific tuning---a hidden cost that is
not visible from synthetic evaluation alone.

These findings establish four conclusions about agent failure detection.
First, it is a genuinely difficult problem: open-weight LLM judges that
achieve perfect detection on synthetic data fail badly on real traces, and
even frontier models (Claude Sonnet~4) achieve category alignment below 0.5.
Second, the failure modes of heuristic and LLM-based systems are largely
complementary: heuristics catch structural failures with zero false positives,
and LLM judges catch semantic failures that heuristics cannot see. Claude
Sonnet~4 represents the most capable single-system option, while a
Galea+Claude ensemble represents the most promising direction for
comprehensive agent monitoring. Third, configuration effort is a first-class
evaluation dimension: zero-config approaches (Galea heuristic, LLM judges)
generalize across failure categories without per-deployment authorship, while
rule-driven platforms (LangSmith, Braintrust) scale coverage only with
continued engineering investment---and since both platforms produce identical
DR when given identical rules, the bottleneck is the rules, not the platform.
Fourth, domain calibration is a hidden cost that does not appear in synthetic
evaluation: Braintrust's event-count threshold achieves 91.7\% DR on retail
but 71.1\% FPR, demonstrating that a rule set tuned for one distribution can
be operationally unusable on a different domain without further calibration.

\paragraph{Future work.}
This initial release establishes the benchmark framework and baseline
results. Immediate next steps include: (1)~evaluating GPT-4.1 as an
additional proprietary LLM judge to complete the planned system comparison,
and evaluating LangSmith and Braintrust each with a more comprehensive
rule set (covering all 7 failure categories) to measure the engineering
investment required to close the coverage gap;
(2)~evaluating the Galea full-stack system (with LLM-augmented analysis)
beyond the heuristic-only configuration reported here;
(3)~acquiring and integrating \swebench trajectories;
(4)~running the multi-judge consensus pipeline for ground truth
validation; (5)~applying domain-specific quality rubrics to enable the
\texttt{degraded} tier and QDR measurement; (6)~expanding domain coverage
to medical, legal, and financial workflows; (7)~adding streaming
evaluation for real-time monitoring tools; (8)~investigating why
category alignment is so low and whether alternative prompting strategies
or fine-tuned models can improve failure type identification; and
(9)~formally evaluating heuristic+LLM ensemble configurations to test the
complementarity hypothesis.

% ==========================================================================
% REFERENCES
% ==========================================================================
\bibliography{atfd}

% ==========================================================================
% APPENDIX
% ==========================================================================
\appendix
\newpage

\section{Example Trajectory}
\label{app:trajectory}

Figure~\ref{fig:example_trajectory} shows an annotated example trajectory
from the \taubench retail domain with injected \texttt{action.wrong\_args}
and \texttt{communication.missing\_info} failures.

\begin{figure}[h]
\centering
\begin{lstlisting}[caption={Example trajectory with two failures: the
exchange tool is called with the wrong target item (action.wrong\_args)
and the confirmation message omits the refund amount
(communication.missing\_info).}, label={fig:example_trajectory}]
Event 1 [user_message]:
  "I'd like to exchange the blue shirt (item #1234)
   for a red one (item #5678)."

Event 2 [tool_call]:
  exchange_item(order_id="ORD-001",
                source_item="1234",
                target_item="9999")     # <-- WRONG_ARGS: should be 5678

Event 3 [tool_result]:
  {"status": "success",
   "refund_amount": 12.50,
   "new_item": "9999"}

Event 4 [assistant_message]:
  "I've processed the exchange for you.
   Your new item will arrive in 3-5
   business days."                      # <-- MISSING_INFO: no refund amount
\end{lstlisting}
\end{figure}

\section{Full Confusion Matrix Template}
\label{app:confusion}

Table~\ref{tab:confusion} shows the confusion matrix template used for
per-system analysis. Each cell counts the number of trajectories where the
ground-truth outcome (rows) was mapped to the predicted outcome (columns).

\begin{table}[h]
\centering
\caption{Confusion matrix template for three-tier outcome classification.}
\label{tab:confusion}
\small
\begin{tabular}{@{}lccc@{}}
\toprule
& \multicolumn{3}{c}{\textbf{Predicted}} \\
\cmidrule(l){2-4}
\textbf{Ground Truth} & Pass & Degraded & Fail \\
\midrule
Pass     & --- & --- & --- \\
Degraded & \multicolumn{3}{c}{\emph{not yet labeled}} \\
Fail     & --- & --- & --- \\
\bottomrule
\end{tabular}
\end{table}

\section{Judge Prompt Template}
\label{app:judge_prompt}

The following is the two-stage prompt template used for the LLM judge
systems. Stage~1 elicits an unconstrained analysis; Stage~2 maps findings
to the \atfd taxonomy.

\paragraph{Stage 1: Open Analysis.}

\begin{lstlisting}[language={},breaklines=true]
You are an expert agent trajectory analyst. You will be
given a complete agent trajectory consisting of user
messages, assistant messages, tool calls, and tool results.

Your task is to carefully analyze the trajectory and
identify ALL potential issues, failures, or quality
concerns. For each issue found, describe:
1. What went wrong
2. Where in the trajectory it occurred (step number)
3. How severe the issue is
4. What the correct behavior should have been

Be thorough. Consider action correctness, state
consistency, communication quality, process efficiency,
safety compliance, and output quality.

TRAJECTORY:
{trajectory_json}

TASK DESCRIPTION:
{task_description}
\end{lstlisting}

\paragraph{Stage 2: Taxonomy Classification.}

\begin{lstlisting}[language={},breaklines=true]
Given your analysis above, now classify each identified
issue using the ATFD failure taxonomy.

For each finding, output a JSON object with:
- "severity": "error" | "warning" | "info"
- "category": "<category>.<subcategory>" from the
  taxonomy below
- "description": one-sentence explanation
- "attribution": step number or range where the
  failure occurred

TAXONOMY:
{taxonomy_json}

Also set "has_failure" to true if ANY error or warning
finding exists, false otherwise.

Output valid JSON matching this schema:
{output_schema}
\end{lstlisting}

\end{document}
